\documentclass[aspectratio=169]{beamer}

% Tema UFS — Universidade Federal de Sergipe
\usetheme{UFS}

% Pacotes base
\usepackage[utf8]{inputenc}
\usepackage[portuguese]{babel}
\usepackage{amsmath,amssymb}
\usepackage{graphicx}
\usepackage{booktabs}
\usepackage{xcolor}

% TikZ e bibliotecas
\usepackage{tikz}
\usetikzlibrary{shapes.geometric, arrows.meta, positioning, matrix, fit,
  backgrounds, decorations.pathreplacing, shadows, patterns}

% Listagens — paleta VS Code Light+
\usepackage{listings}
\definecolor{bgcolor}{HTML}{FFFFFF}
\definecolor{linecolor}{HTML}{DDDDDD}
\definecolor{numcolor}{HTML}{999999}
\definecolor{kwcolor}{HTML}{0000FF}
\definecolor{strcolor}{HTML}{A31515}
\definecolor{cmtcolor}{HTML}{008000}

\lstdefinestyle{ufsStyle}{
  backgroundcolor=\color{bgcolor},
  basicstyle=\ttfamily\footnotesize\color{black},
  keywordstyle=\color{kwcolor}\bfseries,
  stringstyle=\color{strcolor},
  commentstyle=\color{cmtcolor}\itshape,
  numberstyle=\tiny\color{numcolor},
  numbers=left, numbersep=12pt,
  xleftmargin=20pt, framexleftmargin=20pt,
  breaklines=true, tabsize=4,
  keepspaces=true, showspaces=false,
  showstringspaces=false, showtabs=false,
  frame=single, rulecolor=\color{linecolor},
  framesep=4pt, captionpos=b, upquote=true,
  columns=fullflexible,
}
\lstset{style=ufsStyle, language=C}

% --- Metadados ---
\title{Programação C}
\subtitle{Tipos de dados, estruturas de dados e tipos abstratos de dados}
\author{Prof.\ André Yoshiaki Kashiwabara}
\institute{DCOMP -- Universidade Federal de Sergipe\\
\texttt{andreyoshiaki@dcomp.ufs.br}}
\date{}

\begin{document}

\begin{frame}[plain]
  \titlepage
\end{frame}

\begin{frame}{Agenda}
  \tableofcontents
\end{frame}

\begin{frame}{O que você vai aprender hoje}
  \begin{center}
  \begin{tikzpicture}[
    item/.style={rectangle, draw=UFSBlue, fill=UFSBlue!10, rounded corners,
      minimum width=10.5cm, minimum height=0.8cm, font=\small}
  ]
    \node[item] (t1) at (0,0)     {O que é um \textbf{tipo de dados}: como interpretar os bits e o que fazer com eles};
    \node[item] (t2) at (0,-1.2)  {O que é uma \textbf{estrutura de dados}: como organizar muitos valores na memória};
    \node[item] (t3) at (0,-2.4)  {O que é um \textbf{tipo abstrato de dados}: o que se pode fazer, sem dizer como};
    \draw[->, thick, UFSBlue] (t1)--(t2);
    \draw[->, thick, UFSBlue] (t2)--(t3);
  \end{tikzpicture}
  \end{center}
\end{frame}

% =====================================================================
\section{O programa de hoje}

\begin{frame}[shrink=5]{De onde partimos}
  \begin{columns}[T]
    \begin{column}{0.5\textwidth}
      \textbf{Na aula passada:}\\[4pt]
      Um programa que lê 5 notas num vetor, calcula a média numa função e
      classifica a turma.
    \end{column}
    \begin{column}{0.5\textwidth}
      \textbf{Hoje:}\\[4pt]
      O mesmo problema cresceu: as notas ficam \textbf{cadastradas} e o
      programa oferece \textbf{operações} sobre a coleção.
    \end{column}
  \end{columns}
  \vspace{1em}
  \begin{center}
  \begin{tikzpicture}[
    lente/.style={rectangle, draw=UFSBlue, fill=UFSBlue!10, rounded corners,
      minimum width=3.6cm, minimum height=0.9cm, font=\small, align=center},
    prog/.style={rectangle, draw=UFSGray, fill=UFSGray!10, rounded corners,
      minimum width=3.2cm, minimum height=0.9cm, font=\small\ttfamily}
  ]
    \node[prog] (p) at (0,0) {notas\_turma.c};
    \node[lente] (l1) at (-4.6,-2) {Tipos de\\dados};
    \node[lente] (l2) at (0,-2)    {Estruturas de\\dados};
    \node[lente] (l3) at (4.6,-2)  {Tipos abstratos\\de dados};
    \draw[->, thick, UFSBlue] (p) -- (l1);
    \draw[->, thick, UFSBlue] (p) -- (l2);
    \draw[->, thick, UFSBlue] (p) -- (l3);
  \end{tikzpicture}
  \end{center}
  \begin{center}
    Vamos olhar um único programa por \textbf{três lentes} diferentes.
  \end{center}
\end{frame}

\begin{frame}[fragile,shrink=5]{O programa: \texttt{inserir} e \texttt{media}}
\begin{lstlisting}
#include <stdio.h>
#define CAPACIDADE 8

int inserir(double v[], int n, double nota) {
    if (n == CAPACIDADE)
        return n;          /* colecao cheia: ignora a nota */
    v[n] = nota;
    return n + 1;
}

double media(double v[], int n) {
    double soma = 0.0;
    for (int i = 0; i < n; i++)
        soma += v[i];
    return soma / n;
}
\end{lstlisting}
\end{frame}

\begin{frame}[fragile]{O programa: \texttt{maior}}
\begin{lstlisting}
double maior(double v[], int n) {
    double m = v[0];
    for (int i = 1; i < n; i++)
        if (v[i] > m)
            m = v[i];
    return m;
}
\end{lstlisting}
  \begin{itemize}
    \item Repare no padrão: toda operação recebe o \textbf{vetor} e o
      \textbf{contador} \texttt{n} --- e mais nada.
    \item Leia com calma: você vai prever a saída do programa completo no
      próximo passo.
  \end{itemize}
\end{frame}

\begin{frame}[fragile,shrink=5]{O programa: o \texttt{main}}
\begin{lstlisting}
int main(void) {
    double notas[CAPACIDADE];
    int n = 0;

    n = inserir(notas, n, 7.0);
    n = inserir(notas, n, 8.5);
    n = inserir(notas, n, 6.0);

    printf("%d notas cadastradas\n", n);
    printf("media: %.2f\n", media(notas, n));
    printf("maior: %.1f\n", maior(notas, n));
    printf("cada nota ocupa %zu bytes\n", sizeof(double));
    printf("o vetor ocupa %zu bytes\n", sizeof notas);
    return 0;
}
\end{lstlisting}
\end{frame}

\begin{frame}{O que este programa imprime?}
  \begin{block}{Antes de executar, preveja --- em dupla, 2--3 minutos}
    \begin{enumerate}
      \item Escreva as \textbf{cinco linhas} de saída, exatamente como vão
        aparecer no terminal.
      \item Que valor sai como média? Com quantas casas decimais?
      \item O vetor guarda \textbf{3 notas} --- quantos bytes você acha que a
        última linha vai informar? Por quê?
    \end{enumerate}
  \end{block}
  \vspace{0.5em}
  Anote as previsões: vamos compará-las com a saída real já já.
\end{frame}

\begin{frame}[fragile]{Compile e execute}
\begin{lstlisting}[language=bash, numbers=none]
$ gcc notas_turma.c -o notas_turma
$ ./notas_turma
3 notas cadastradas
media: 7.17
maior: 8.5
cada nota ocupa 8 bytes
o vetor ocupa 64 bytes
\end{lstlisting}
  \begin{itemize}
    \item Compare com a sua previsão: o que bateu? O que surpreendeu?
    \item Duas surpresas comuns: a média \texttt{7.17} (e não \texttt{7.16}?)
      e os \texttt{64 bytes} (e não 24?).
    \item As respostas passam pelos três conceitos de hoje --- vamos a eles.
  \end{itemize}
\end{frame}

% =====================================================================
\section{Tipos de dados}

\begin{frame}{Por que precisamos de tipos?}
  \begin{columns}[T]
    \begin{column}{0.5\textwidth}
      \textbf{Problema:}\\[4pt]
      Para o computador, a memória é só uma sequência de \textbf{bits}.
      O padrão \texttt{01000001} é a letra \texttt{A}? O número 65? Parte de
      um número maior?\\[4pt]
      Os bits, sozinhos, \textbf{não dizem o que significam}.
    \end{column}
    \begin{column}{0.5\textwidth}
      \textbf{Solução:}\\[4pt]
      O \textbf{tipo de dados} responde as duas perguntas que os bits não
      respondem:
      \begin{itemize}
        \item \textbf{como interpretar} os bits;
        \item \textbf{o que se pode fazer} com eles.
      \end{itemize}
    \end{column}
  \end{columns}
\end{frame}

\begin{frame}{Tipo de dados: definição}
  \begin{block}{Definição}
    Um \textbf{tipo de dados} é um conjunto de \textbf{valores} possíveis,
    as \textbf{operações} permitidas sobre eles e a \textbf{representação}
    desses valores na memória (quantos bytes, em que formato).
  \end{block}
  \begin{exampleblock}{Intuição}
    O tipo é um \textbf{contrato de leitura}: os mesmos bits, lidos com
    contratos diferentes, significam coisas diferentes --- como uma sequência
    de letras que muda de sentido conforme o idioma em que você a lê.
  \end{exampleblock}
\end{frame}

\begin{frame}{Os mesmos bits, duas leituras}
  \begin{center}
  \begin{tikzpicture}[
    bitcell/.style={rectangle, draw=UFSBlue, fill=UFSBlue!10,
      minimum width=0.62cm, minimum height=0.9cm, font=\ttfamily\small},
    leitura/.style={rectangle, draw=UFSGray, fill=UFSGray!10, rounded corners,
      minimum width=4.2cm, minimum height=0.9cm, font=\small, align=center}
  ]
    \foreach \i/\b in {0/0, 1/1, 2/0, 3/0, 4/0, 5/0, 6/0, 7/1} {
      \node[bitcell] (b\i) at (\i*0.62,0) {\b};
    }
    \node[font=\scriptsize, text=UFSGray] at (2.2,0.85) {um byte na memória};
    \node[leitura] (c) at (-2.2,-2) {lido como \texttt{char}:\\o caractere \texttt{'A'}};
    \node[leitura] (i) at (6.6,-2)  {lido como número:\\o inteiro \texttt{65}};
    \draw[->, thick, UFSBlue] (b1.south) -- (c.north);
    \draw[->, thick, UFSBlue] (b6.south) -- (i.north);
  \end{tikzpicture}
  \end{center}
  \begin{itemize}
    \item O tipo também decide a \textbf{operação}: \texttt{7 / 2} vale
      \texttt{3} (divisão inteira), mas \texttt{7.0 / 2} vale \texttt{3.5}.
    \item Na aula passada: trocar \texttt{soma} para \texttt{int} mudou a
      média --- mesmo programa, outro contrato de leitura.
  \end{itemize}
\end{frame}

\begin{frame}[fragile]{Os tipos básicos de C}
  \begin{center}
  \begin{tabular}{llll}
    \toprule
    Tipo & Guarda & Tamanho típico & Exemplo de literal \\
    \midrule
    \texttt{char}   & um caractere        & 1 byte  & \texttt{'A'} \\
    \texttt{int}    & inteiro             & 4 bytes & \texttt{42} \\
    \texttt{float}  & real, $\sim$7 dígitos  & 4 bytes & \texttt{3.14f} \\
    \texttt{double} & real, $\sim$15 dígitos & 8 bytes & \texttt{3.14} \\
    \bottomrule
  \end{tabular}
  \end{center}
  \vspace{0.5em}
  \begin{itemize}
    \item O operador \texttt{sizeof} informa quantos \textbf{bytes} um tipo
      (ou uma variável) ocupa --- o resultado tem tipo \texttt{size\_t} e é
      impresso com \texttt{\%zu}.
    \item ``Típico'' porque o padrão C fixa \textbf{mínimos}, não valores
      exatos: os tamanhos podem variar entre máquinas.
  \end{itemize}
\end{frame}

\begin{frame}{A lente dos tipos sobre o programa}
  \begin{itemize}
    \item \texttt{cada nota ocupa 8 bytes} --- é o \texttt{sizeof(double)}:
      a \textbf{representação} faz parte do tipo.
    \item A média exata é $21{,}5 / 3 = 7{,}1\overline{6}$; o
      \texttt{double} guarda o valor com $\sim$15 dígitos e o \texttt{\%.2f}
      \textbf{arredonda} para \texttt{7.17} --- não trunca.
    \item Se \texttt{soma} fosse \texttt{int}, a divisão \texttt{soma / n}
      seria inteira: média \texttt{7.00}. O \textbf{tipo escolhe a operação}.
    \item \texttt{n} é \texttt{int} porque conta coisas; \texttt{notas}
      guarda \texttt{double} porque notas têm casas decimais. Escolher tipo é
      escolher \textbf{o que o valor significa}.
  \end{itemize}
\end{frame}

\begin{frame}{Recapitulando: tipos de dados}
  \begin{block}{O que aprendemos}
    \begin{itemize}
      \item Bits não têm significado próprio; o \textbf{tipo} diz como
        interpretá-los.
      \item Tipo = \textbf{valores} + \textbf{operações} +
        \textbf{representação}.
      \item \texttt{sizeof} revela a representação; a mesma expressão muda de
        resultado conforme o tipo (\texttt{7/2} vs.\ \texttt{7.0/2}).
    \end{itemize}
  \end{block}
  \begin{exampleblock}{Pergunta-guia deste conceito}
    \emph{``Como interpreto um valor e o que posso fazer com ele?''}
  \end{exampleblock}
\end{frame}

% =====================================================================
\section{Estruturas de dados}

\begin{frame}{Por que precisamos de estruturas de dados?}
  \begin{columns}[T]
    \begin{column}{0.5\textwidth}
      \textbf{Problema:}\\[4pt]
      Uma variável de um tipo básico guarda \textbf{um} valor. A turma tem
      dezenas de notas; um texto tem milhares de caracteres.\\[4pt]
      Criar \texttt{nota1}, \texttt{nota2}, \dots, \texttt{nota50} não
      escala.
    \end{column}
    \begin{column}{0.5\textwidth}
      \textbf{Solução:}\\[4pt]
      Uma \textbf{estrutura de dados}: uma forma de \textbf{organizar} muitos
      valores na memória, com \textbf{regras de acesso} conhecidas.\\[4pt]
      O vetor é a primeira que você já conhece.
    \end{column}
  \end{columns}
\end{frame}

\begin{frame}{Estrutura de dados: definição}
  \begin{block}{Definição}
    Uma \textbf{estrutura de dados} é uma forma de organizar e relacionar
    dados na memória, junto com as \textbf{regras de acesso} a esses dados
    (onde inserir, como percorrer, como localizar).
  \end{block}
  \begin{exampleblock}{Intuição}
    Uma biblioteca não é só uma pilha de livros: é a estante, a ordem nas
    prateleiras e o catálogo. A \textbf{organização} é o que torna o acesso
    rápido --- os dados são os mesmos, o arranjo é que muda tudo.
  \end{exampleblock}
\end{frame}

\begin{frame}[shrink=8]{A estrutura do programa: vetor + contador}
  \begin{center}
  \begin{tikzpicture}[
    cel/.style={rectangle, draw=UFSBlue, fill=UFSBlue!10,
      minimum width=1.35cm, minimum height=0.9cm, font=\small},
    vazia/.style={rectangle, draw=UFSGray, fill=white,
      minimum width=1.35cm, minimum height=0.9cm, font=\scriptsize,
      text=UFSGray},
    idx/.style={font=\scriptsize, text=UFSGray}
  ]
    \node[cel] (c0) at (0,0)      {7.0};
    \node[cel] (c1) at (1.35,0)   {8.5};
    \node[cel] (c2) at (2.7,0)    {6.0};
    \foreach \i in {3,...,7} {
      \node[vazia] (c\i) at (\i*1.35,0) {?};
    }
    \foreach \i in {0,...,7} {
      \node[idx] at (\i*1.35,0.8) {[\i]};
    }
    \draw[decorate, decoration={brace, mirror, amplitude=6pt}, UFSBlue]
      (-0.67,-0.75) -- (3.37,-0.75) node[midway, below=8pt, font=\scriptsize]
      {\texttt{n = 3} posições usadas};
    \draw[decorate, decoration={brace, mirror, amplitude=6pt}, UFSGray]
      (-0.67,-1.7) -- (10.12,-1.7) node[midway, below=8pt, font=\scriptsize]
      {\texttt{CAPACIDADE = 8} posições reservadas --- $8 \times 8 = 64$ bytes};
  \end{tikzpicture}
  \end{center}
  \begin{itemize}
    \item Mistério dos 64 bytes resolvido: \texttt{sizeof} mede a
      \textbf{capacidade reservada}, não o quanto está em uso.
    \item A estrutura de hoje é a \textbf{dupla} vetor + \texttt{n}: os dados,
      a organização (posições contíguas) e a regra (as \texttt{n} primeiras
      valem).
    \item C não verifica limites: a regra \texttt{n <= CAPACIDADE} é
      \textbf{nossa} responsabilidade --- por isso o teste em \texttt{inserir}.
  \end{itemize}
\end{frame}

\begin{frame}{Nem tudo é vetor: o panorama do semestre}
  \begin{center}
  \begin{tikzpicture}[
    cel/.style={rectangle, draw=UFSBlue, fill=UFSBlue!10,
      minimum width=0.7cm, minimum height=0.55cm, font=\scriptsize},
    no/.style={rectangle, draw=UFSBlue, fill=UFSBlue!10, rounded corners=2pt,
      minimum width=0.7cm, minimum height=0.55cm, font=\scriptsize},
    rot/.style={font=\small\bfseries, text=UFSBlue},
    leg/.style={font=\scriptsize, text=UFSGray, align=center}
  ]
    % Vetor
    \node[rot] at (0.7,1.1) {Vetor};
    \foreach \i in {0,1,2,3} \node[cel] at (\i*0.7,0) {};
    \node[leg] at (1.05,-0.8) {acesso direto\\pelo índice};
    % Lista encadeada
    \node[rot] at (5.3,1.1) {Lista encadeada};
    \node[no] (n1) at (4.2,0) {};
    \node[no] (n2) at (5.5,0) {};
    \node[no] (n3) at (6.8,0) {};
    \draw[->, UFSBlue] (n1) -- (n2);
    \draw[->, UFSBlue] (n2) -- (n3);
    \node[leg] at (5.5,-0.8) {cresce e encolhe\\sem capacidade fixa};
    % Pilha
    \node[rot] at (9.3,1.1) {Pilha};
    \foreach \i in {0,1,2} \node[cel] at (9.3,-0.35+\i*0.55) {};
    \draw[->, UFSBlue] (10.3,1.0) -- (9.75,0.75);
    \node[leg] at (9.3,-1.2) {entra e sai\\pelo topo};
    % Fila
    \node[rot] at (13,1.1) {Fila};
    \foreach \i in {0,1,2} \node[cel] at (12.3+\i*0.7,0) {};
    \draw[->, UFSBlue] (11.3,0) -- (11.85,0);
    \draw[->, UFSBlue] (13.45,0) -- (14,0);
    \node[leg] at (13,-0.8) {entra no fim,\\sai da frente};
  \end{tikzpicture}
  \end{center}
  \vspace{0.5em}
  \begin{itemize}
    \item Cada organização favorece um uso: não existe estrutura
      ``melhor'' --- existe a \textbf{adequada ao problema}.
    \item Neste semestre vamos \textbf{implementar} cada uma delas em C, com
      \texttt{struct} e ponteiros.
  \end{itemize}
\end{frame}

\begin{frame}{Recapitulando: estruturas de dados}
  \begin{block}{O que aprendemos}
    \begin{itemize}
      \item Estrutura de dados = \textbf{organização} + \textbf{regras de
        acesso}, não só os dados.
      \item Vetor + contador \texttt{n}: contíguo, capacidade fixa, as
        \texttt{n} primeiras posições valem.
      \item Manter as regras (ex.\ não passar da capacidade) é
        responsabilidade do programa em C.
    \end{itemize}
  \end{block}
  \begin{exampleblock}{Pergunta-guia deste conceito}
    \emph{``Como organizo muitos valores para acessá-los do jeito que o
    problema pede?''}
  \end{exampleblock}
\end{frame}

% =====================================================================
\section{Tipos abstratos de dados}

\begin{frame}{Por que precisamos de abstração?}
  \begin{columns}[T]
    \begin{column}{0.5\textwidth}
      \textbf{Problema:}\\[4pt]
      Se quem \textbf{usa} a coleção de notas precisa conhecer o vetor, o
      \texttt{n} e a capacidade, qualquer mudança interna quebra todo o
      programa --- e cada usuário pode violar as regras sem querer.
    \end{column}
    \begin{column}{0.5\textwidth}
      \textbf{Solução:}\\[4pt]
      Separar \textbf{o quê} se pode fazer (as operações e suas regras)
      de \textbf{como} é feito (a estrutura interna).\\[4pt]
      Essa separação tem nome: \textbf{tipo abstrato de dados} (TAD).
    \end{column}
  \end{columns}
\end{frame}

\begin{frame}{Tipo abstrato de dados: definição}
  \begin{block}{Definição}
    Um \textbf{tipo abstrato de dados} é a especificação de um conjunto de
    \textbf{operações} sobre uma coleção de valores e das \textbf{regras de
    comportamento} dessas operações --- \textbf{sem fixar} como os dados são
    representados nem como as operações são implementadas.
  \end{block}
  \begin{exampleblock}{Intuição}
    Um caixa eletrônico: você conhece as operações (consultar saldo, sacar,
    depositar) e as regras (não sacar mais que o saldo). Como o banco guarda
    as contas por dentro? Você não sabe --- e \textbf{não precisa} saber.
  \end{exampleblock}
\end{frame}

\begin{frame}[shrink=10]{A fronteira entre o quê e o como}
  \begin{center}
  \begin{tikzpicture}[
    usuario/.style={rectangle, draw=UFSGray, fill=UFSGray!10, rounded corners,
      minimum width=2.6cm, minimum height=0.9cm, font=\small\ttfamily},
    op/.style={rectangle, draw=UFSBlue, fill=UFSBlue!10, rounded corners,
      minimum width=2.6cm, minimum height=0.75cm, font=\small\ttfamily},
    interno/.style={rectangle, draw=UFSBlue!60, fill=white, dashed,
      minimum width=3.4cm, minimum height=0.75cm, font=\small\ttfamily}
  ]
    \node[usuario] (main) at (0,0) {main()};
    \node[op] (op1) at (5,1.3)  {inserir()};
    \node[op] (op2) at (5,0)    {media()};
    \node[op] (op3) at (5,-1.3) {maior()};
    \node[interno] (rep) at (10.2,0) {double v[8]; int n;};
    \draw[->, thick, UFSBlue] (main) -- (op1);
    \draw[->, thick, UFSBlue] (main) -- (op2);
    \draw[->, thick, UFSBlue] (main) -- (op3);
    \draw[->, UFSBlue!60, dashed] (op1) -- (rep);
    \draw[->, UFSBlue!60, dashed] (op2) -- (rep);
    \draw[->, UFSBlue!60, dashed] (op3) -- (rep);
    \draw[very thick, UFSBlue] (7.3,2) -- (7.3,-2)
      node[below=2pt, font=\scriptsize, text=UFSBlue] {fronteira da abstração};
    \node[font=\scriptsize, text=UFSGray] at (2.5,2) {o quê: operações};
    \node[font=\scriptsize, text=UFSGray] at (10.2,2) {como: representação};
  \end{tikzpicture}
  \end{center}
  \begin{itemize}
    \item O \texttt{main} do nosso programa só cruza a fronteira pelas
      \textbf{operações} --- ele nunca toca \texttt{v[]} diretamente.
    \item A regra ``não passar da capacidade'' vive \textbf{dentro} de
      \texttt{inserir}: quem usa não consegue violá-la.
  \end{itemize}
\end{frame}

\begin{frame}{Você já usa TADs sem saber}
  \begin{itemize}
    \item \texttt{printf("media: \%.2f\textbackslash n", m)} --- você conhece
      \textbf{o que} ele faz; sabe \textbf{como} ele converte o
      \texttt{double} em texto e o entrega à tela? Não precisa.
    \item As funções de \texttt{math.h} (\texttt{sqrt}, \texttt{pow}):
      operações com comportamento conhecido, implementação escondida.
    \item Toda a biblioteca padrão de C é organizada assim: \textbf{contratos}
      no \texttt{.h}, implementação em outro lugar.
  \end{itemize}
  \begin{exampleblock}{A consequência prática}
    Se amanhã a coleção de notas virar uma lista encadeada, só os corpos de
    \texttt{inserir}, \texttt{media} e \texttt{maior} mudam --- o
    \texttt{main} continua \textbf{idêntico}. Abstração é o que permite um
    programa \textbf{evoluir} sem quebrar.
  \end{exampleblock}
\end{frame}

\begin{frame}{Os três conceitos, lado a lado}
  \begin{center}
  \begin{tabular}{p{3.1cm}p{4.6cm}p{4.6cm}}
    \toprule
    Conceito & Pergunta-guia & No programa de hoje \\
    \midrule
    Tipo de dados &
    Como interpreto um valor e o que posso fazer com ele? &
    \texttt{double}, \texttt{int}, \texttt{sizeof}, arredondamento \\[6pt]
    Estrutura de dados &
    Como organizo muitos valores na memória? &
    vetor contíguo + contador \texttt{n} \\[6pt]
    Tipo abstrato de dados &
    O que se pode fazer, sem dizer como? &
    \texttt{inserir}, \texttt{media}, \texttt{maior} escondendo o vetor \\
    \bottomrule
  \end{tabular}
  \end{center}
  \vspace{0.5em}
  \begin{center}
    Três níveis da mesma escada: \textbf{interpretar} um valor,
    \textbf{organizar} muitos, \textbf{esconder} a organização.
  \end{center}
\end{frame}

\begin{frame}{Recapitulando: tipos abstratos de dados}
  \begin{block}{O que aprendemos}
    \begin{itemize}
      \item TAD = \textbf{operações} + \textbf{regras de comportamento}, sem
        fixar a implementação.
      \item A fronteira da abstração protege as regras e isola quem usa de
        quem implementa.
      \item Em C, a fronteira se materializa em \textbf{funções} (e, adiante,
        em \texttt{struct} + headers).
    \end{itemize}
  \end{block}
  \begin{exampleblock}{Pergunta-guia deste conceito}
    \emph{``O que esta coleção faz --- independente de como é feita por
    dentro?''}
  \end{exampleblock}
\end{frame}

% =====================================================================
\section{Mãos à obra}

\begin{frame}[shrink=8]{Altere o programa}
  \begin{block}{Tarefas, em ordem crescente de desafio --- uma por lente}
    \begin{enumerate}
      \item \textbf{Tipos:} troque \texttt{double} por \texttt{float} em todo
        o programa. O que muda nas duas últimas linhas da saída? E na média?
      \item \textbf{Estruturas:} num laço, insira 10 notas iguais a
        \texttt{5.0}. Quantas entram? Como você descobre, pelo programa, que
        as demais foram ignoradas?
      \item \textbf{Abstração:} acrescente a operação
        \texttt{double menor(double v[], int n)} e imprima também a menor
        nota --- sem alterar nenhuma das operações existentes.
    \end{enumerate}
  \end{block}
  \begin{exampleblock}{Dica}
    Para cada tarefa: altere, recompile e \textbf{preveja a saída antes de
    executar}.
  \end{exampleblock}
\end{frame}

\begin{frame}[shrink=8]{Desafio: registro de temperaturas}
  \begin{block}{Desafio}
    Uma estação meteorológica registra leituras de temperatura ao longo do
    dia. Escreva as operações de um registro de temperaturas:
    \begin{itemize}
      \item \texttt{int registrar(double v[], int n, double t)} --- aceita
        somente leituras entre $-10$ e $50\,^{\circ}$C (regra dentro da
        operação!);
      \item \texttt{double media(double v[], int n)};
      \item \texttt{double amplitude(double v[], int n)} --- maior menos
        menor leitura.
    \end{itemize}
    O \texttt{main} deve usar \textbf{apenas} essas operações.
  \end{block}
  \begin{exampleblock}{Terminou? Estenda}
    Conte quantas leituras foram \textbf{rejeitadas} por estarem fora da
    faixa --- sem quebrar a fronteira da abstração.
  \end{exampleblock}
\end{frame}

\begin{frame}{Recapitulando}
  \begin{block}{Os três conceitos de hoje}
    \begin{itemize}
      \item \textbf{Tipo de dados:} valores + operações + representação ---
        o contrato de leitura dos bits.
      \item \textbf{Estrutura de dados:} organização + regras de acesso ---
        o arranjo dos valores na memória.
      \item \textbf{Tipo abstrato de dados:} operações + comportamento ---
        o quê, sem o como.
    \end{itemize}
  \end{block}
  \begin{exampleblock}{Como trabalhamos}
    Prevendo, executando, investigando, modificando e criando --- como será
    em todo o semestre.
  \end{exampleblock}
\end{frame}

\begin{frame}{Para casa}
  \begin{itemize}
    \item \textbf{Próxima aula (24/08):} registros (\texttt{struct}) ---
      declaração, acesso, structs aninhadas, alinhamento e padding na
      memória. É a primeira estrutura que vamos construir de verdade.
    \item Fazer o tutorial guiado (Jupyter Notebook) do Classroom --- ele
      continua as tarefas de alteração e o desafio para quem não terminou em
      sala.
    \item Ler o capítulo introdutório de Tenenbaum, Langsam \& Augenstein
      (\emph{Estruturas de dados usando C}) sobre tipos e abstração.
    \item Dúvidas: \texttt{andreyoshiaki@dcomp.ufs.br} ou no início da
      próxima aula.
  \end{itemize}
\end{frame}

% =====================================================================
\section*{Referências}
\begin{frame}[allowframebreaks]{Referências}
  \footnotesize
  \nocite{*}
  \bibliographystyle{plain}
  \bibliography{../referencias}
\end{frame}

\end{document}
